%=============================================================================
%  Physical Intelligence pi0.5 on the Unified Engine v1
%  Single-column report.
%  Build: pdflatex main.tex   (article, amsmath, booktabs)
%  Overleaf: paste as main.tex, compiler = pdfLaTeX.
%=============================================================================
\documentclass[11pt]{article}

\usepackage[letterpaper,margin=1in]{geometry}
\usepackage{amsmath,amssymb}
\usepackage{booktabs}
\usepackage{float}
\usepackage{xcolor}
\usepackage{graphicx}
\usepackage[hidelinks]{hyperref}

% ---- CONFIDENTIAL watermark on every page -------------------------------
\AddToHook{shipout/background}{%
  \put(0.5\paperwidth,-0.5\paperheight){%
    \makebox(0,0){\rotatebox{45}{%
      \textcolor{gray!25}{\fontsize{68}{68}\selectfont\bfseries CONFIDENTIAL}}}}%
}

% ---- macros -------------------------------------------------------------
\newcommand{\pizero}{$\pi_{0.5}$}
\newcommand{\ue}{Unified~Engine}
\newcommand{\iffour}{IF4}
\newcommand{\bfteen}{bf16}
\newcommand{\numTBD}{{\color{red}--}}

\emergencystretch=1.5em              % absorb marginal overfulls in justification
\setlength{\textfloatsep}{6pt}       % tighten space around floats
\setlength{\floatsep}{6pt}
\setlength{\intextsep}{6pt}
\setlength{\parskip}{2pt}

\title{\vspace{-2em}Deploying Physical Intelligence \pizero{} on the \ue{}}
\author{Apex Compute}
\date{\vspace{-1em}}

\begin{document}
\maketitle

%=============================================================================
\section{The \pizero{} Model}
\label{sec:model}
%=============================================================================
\pizero{} is a vision--language--action (VLA) policy of roughly three billion
parameters that maps camera observations and a language instruction directly to
robot actions. Each inference emits a $10{\times}7$ action chunk: ten future
control steps, each a $7$-dimensional command giving the robot's degrees of
freedom (arm motion plus gripper). The computation splits into three stages run
in sequence---a vision encoder, a language prefix, and a flow-matching action
expert.

\subsection{Vision encoder}
Each image is encoded independently by a SigLIP ViT; \pizero{} takes up to three
slots (scene, wrist, one padded). Every $224{\times}224$ image becomes $256$
patch tokens of width $1152$, run through $L_v{=}27$ bidirectional layers
($16$ heads, head-dim $72$, MLP width $4304$). Re-run per image and never
cached, it costs $\approx$656\,GFLOP.

\subsection{Language prefix}
The patch tokens are projected to $d{=}2048$, concatenated with the instruction
tokens, and processed once by a Gemma-2B backbone: $L_p{=}18$ layers of
grouped-query attention ($8$ query, $1$ KV head, head-dim $256$), RMSNorm, RoPE,
and a gated SiLU MLP ($I{=}16384$). At prefill length $M_p{=}832$ one layer costs
\begin{equation}
2 M_p d\, d_{q} + 4 M_p d\, d_{kv} + 2 M_p d_{q} d
+ 4 M_p^2 d_{q} + 6 M_p d\, I ,
\end{equation}
($d_q{=}2048$, $d_{kv}{=}256$), dominated by the MLP; over $18$ layers the
prefix is $\approx$3400\,GFLOP ($82\%$ of the model). Each layer's K/V is cached
and reused by every denoise step.

\subsection{Action expert}
A second, smaller Gemma ($\sim$300\,M; $H{=}1024$, MLP $4096$, $18$ layers)
turns noise into the action chunk by flow matching: from Gaussian noise it
integrates a learned velocity field $v_\theta$ over $N{=}10$ Euler steps
($\Delta t{=}-0.1$),
\begin{equation}
x_{t+\Delta t}=x_t+\Delta t\,v_\theta\!\left(x_t,\,t,\,\mathrm{KV}_{\text{prefix}}\right),
\label{eq:euler}
\end{equation}
each step attending over the cached prefix K/V and reading out actions through a
$(1024{\to}32)$ head. At only $\approx$75\,GFLOP ($1.8\%$) it is arithmetically
small but the sole sequential stage, and its per-step error accumulates through
the integrator.

\begin{table}[H]
\centering\small
\caption{Per-stage parameter and FLOP budget (per inference). \emph{Effective}
uses true model dims; \emph{FPGA} is the padded shape the transformer engine
actually issues. The $832$-token prefix is 64-aligned (no padding); the vision
tower pads its attention head-dim ($72{\to}128$) and MLP ($4304{\to}4352$),
$\sim$30\% extra issued work; the action expert
pads its query horizon $10{\to}64$ and action width $32{\to}64$, inflating its
work $\sim$6.5$\times$ --- so overall the FPGA issues $\sim$15\% more FLOPs than
the model nominally needs.}
\label{tab:budget}
\begin{tabular}{@{}lrrr@{}}
\toprule
\textbf{Stage} & \textbf{Params} & \textbf{Eff.\ GFLOP} & \textbf{FPGA GFLOP} \\
\midrule
Vision tower (SigLIP, 3 slots) & $\sim$0.41\,B & 655.9 & 850.7 \\
Language prefix (Gemma-2B)     & $\sim$2.0\,B  & 3399.8 & 3399.8 \\
Action expert (Gemma-300M)     & $\sim$0.31\,B & 74.7 & 483.5 \\
\midrule
\textbf{Total} & \textbf{$\sim$2.7\,B} & \textbf{4130.4} & \textbf{4734.0} \\
\bottomrule
\end{tabular}
\end{table}

%=============================================================================
\section{Porting}
\label{sec:port}
%=============================================================================
The \ue{}~v1 is a soft accelerator IP block whose compute core is the Apex
Compute Transformer Engine; it can be deployed on any
sufficiently large FPGA and operates on $64$-element alignment. In this work a
single transformer engine is targeted to a Xilinx Kintex-7 with $4$\,GB of
attached DRAM. These two
constraints---$64$-alignment and the $4$\,GB DRAM budget---drive the porting
transformations below.


\textbf{\iffour{} weights / DRAM budget.} The $\sim$13\,GB \bfteen{} checkpoint
does not fit the $4$\,GB DRAM budget. All
matmul weights use 4-bit block quantization (block 64, weight-only): every
64-element block stores 64 INT4/FP4 codes plus \emph{one} \bfteen{} scale, so
a quantized matrix costs $4.25$\,bits/weight ($4+16/64$, i.e.\ $3.76\times$
smaller than \bfteen{} rather than $4\times$), while norm gains and biases
stay \bfteen{} outright. That shrinks the three stacks to $\sim$1.56\,GB.
Weights are dequantized to
\bfteen{} immediately before each multiply--accumulate, so \iffour{} is a
\emph{footprint enabler, not a speedup} --- compute is \bfteen{} throughout.

\textbf{Vision head-dim pad.} Head-dim $72$ is not $64$-/$128$-aligned, so the
attention path runs at a padded head-dim of $128$ --- weights pre-padded
$72{\to}128$ across the $Q/K/V/O$ projections and the attention matmuls ---
and is contracted back via a precomputed $128{\times}72$ selection matrix; the
vision MLP pads $4304{\to}4352$. This padding is real issued work, $\sim$30\%
over the effective FLOPs (Table~\ref{tab:budget}).

\textbf{Prefix masking.} When an image slot is unused it can be dropped from
the prefill, compressing $832{\to}576$ (bit-identical, $\sim$24\% faster
end-to-end). Rotary positions use
cumulative-sum-of-mask indices, not \texttt{arange}, so dropping a slot does
not shift phase. The non-causal/causal mixed attention mask is supplied as a
static bias tensor in DRAM. The measured runs in Sec.~\ref{sec:throughput}
keep the full $832$-token prefill so the 3-slot path is exercised end-to-end
(Sec.~\ref{sec:libero}).

\textbf{Denoise loop.} All $10{\times}18$ steps are unrolled into one program
issued with a single execute call; AdaRMSNorm conditioning is a matmul to $3H$
split into scale/shift/gate.

%=============================================================================
\section{Quantization}
\label{sec:quant}
%=============================================================================
\subsection{The Quantization Floor}
\iffour{} is the accelerator's 4-bit weight format: a per-64-block hybrid of
signed INT4 (uniform grid $[-8,7]$) and FP4 E2M1 (non-uniform codebook
$\{0,\pm0.5,\pm1,\pm1.5,\pm2,\pm3,\pm4,\pm6\}$, finer near zero). At pack time
each block is quantized \emph{both} ways and the lower-MSE variant kept; the
choice rides in the \emph{sign} of the block's scale (neg $\to$ INT4, pos $\to$
FP4), so the on-chip dequantizer picks the decode path from one bit, inline
before the MAC---no metadata. INT4 suits dense blocks, FP4 the peaky/outlier
ones. (Results below use the INT4 variant, \texttt{q4\_64}.)

We measure the accuracy cost against a \bfteen{} reference (identical torch run,
same inputs/noise, diffed per stage; Table~\ref{tab:quant},
$\mathrm{SNR}{=}20\log_{10}(\lVert b\rVert/\lVert a{-}b\rVert)$, masked rows
excluded, sample~0)---the theoretical quantization floor, isolated from any
hardware effect. Vision is nearly untouched (25.4\,dB); the prefix caches degrade
with depth, V worse than K ($10.7$ vs.\ $16.0$\,dB mean); and denoise carries the
largest loss as per-step error accumulates through Eq.~\eqref{eq:euler}---$x_t$
walks $45\!\to\!16.3$\,dB (Table~\ref{tab:quant-detail}), concentrating in
low-magnitude action dims (dim\,1, $1.5$\,dB) while task-critical dims hold
(dim\,0 $24.9$, gripper $40.7$\,dB). SNR is a diagnostic; the acceptance gate is
the closed-loop LIBERO agreement (Sec.~\ref{sec:libero}).

\begin{table}[H]
\centering\small
\caption{\iffour{} vs.\ \bfteen{} per section (torch reference, sample~0).}
\label{tab:quant}
\begin{tabular}{@{}lcc@{}}
\toprule
\textbf{Model section} & \textbf{SNR (dB)} & \textbf{Range (dB)} \\
\midrule
Vision encoder (3 cameras)      & 25.37 & 24.3--26.6 \\
Prefix K-cache (18 layers)      & 15.96 & 12.5--19.6 \\
Prefix V-cache (18 layers)      & 10.65 & 7.0--15.4 \\
Denoise velocity $v_t$ (10 steps) & 23.49 & 11.5--26.2 \\
\midrule
\textbf{End-to-end action (10$\times$7)} & \textbf{16.31} & MSE $4.05{\times}10^{-3}$ \\
\bottomrule
\end{tabular}
\end{table}

\begin{table}[H]
\centering\small
\caption{Final-action detail: per-dimension SNR (left) and error accumulation of
the running state $x_t$ across Euler steps (right). Sample~0.}
\label{tab:quant-detail}
\begin{tabular}{@{}lc@{\hspace{1.5em}}lc@{}}
\toprule
\textbf{Action dim} & \textbf{SNR (dB)} & \textbf{Euler step} & \textbf{$x_t$ SNR (dB)} \\
\midrule
dim\,0 (primary) & 24.91 & step 0 & 44.95 \\
dim\,1           & 1.50  & step 3 & 30.13 \\
dim\,2           & 10.08 & step 5 & 24.15 \\
dim\,3           & 6.73  & step 7 & 20.53 \\
dim\,4           & 4.65  & step 8 & 19.26 \\
dim\,5           & 17.61 & step 9 & 16.31 \\
dim\,6 (gripper) & 40.73 & (final) & 16.31 \\
\bottomrule
\end{tabular}
\end{table}

\subsection{LIBERO}
\label{sec:libero}
LIBERO is a standard tabletop manipulation benchmark for evaluating
language-conditioned robot policies: a suite of pick-and-place and articulated
tasks in simulation, scored by closed-loop task success over many episodes. Its
observation uses \textbf{two} camera views---a fixed scene camera and a
wrist-mounted camera. \pizero{}, however, exposes \textbf{three} image slots, and
we process all three: the two real LIBERO cameras plus a third all-zero
(attention-masked) slot. Encoding the empty slot is intentional---other
embodiments and benchmarks drive the policy with a genuine third camera for their
robot actions, so keeping the full 3-slot path exercised on hardware makes the
port drop-in for those settings at no correctness cost (the masked slot
contributes nothing to the output).

\paragraph{An episode is many inferences.} Each model call in
Sec.~\ref{sec:model} predicts a $10{\times}7$ chunk---ten future control
steps. LIBERO runs \emph{closed-loop}: the simulator executes the first
$k{=}10$ steps of that chunk, then feeds the resulting new observation back to
the policy for a fresh chunk, and repeats. So a single task \emph{episode} is a
\emph{sequence} of inferences---one every $k$ sim steps---not a one-shot
prediction, and the episode runs until the task's goal predicate fires (success)
or a per-suite step cap is hit (failure). An episode therefore takes
$\lceil \text{steps}/k \rceil$ inferences: the successful runs in
Table~\ref{tab:libero} are $\sim$$10$--$16$ inferences, the capped failures
$\sim$$23$--$29$. This is what makes the closed-loop test meaningful and also
much heavier than a single forward pass---every one of those inferences runs the
full vision\,+\,prefix\,+\,10-step denoise pipeline on hardware.

Per-tensor SNR is only a diagnostic; the real acceptance test is whether the
quantized policy \emph{acts} the same across that whole loop.
Table~\ref{tab:libero} is the closed-loop comparison---the FPGA \iffour{} policy
vs.\ the CUDA \bfteen{} reference over 10 prompts (9 \texttt{libero\_spatial}
+ 1 \texttt{libero\_object}), one episode each, with identical seed and fixed
denoise noise so both platforms see byte-identical observations. We score each by
task success and by the per-episode step count.

\begin{table}[H]
\centering\small
\caption{LIBERO closed-loop agreement: FPGA (\iffour{}) vs.\ CUDA (\bfteen{}),
10 prompts, 1 episode each, fixed seed/noise. \checkmark${=}$success,
$\times{=}$failure; the number is episode length in sim steps. $\Delta$ is FPGA
$-$ CUDA steps on successes (failures run to the per-suite step cap on both
platforms, so $\Delta$ is
not meaningful there). The two policies reach the \emph{same} outcome on every
prompt.}
\label{tab:libero}
\begin{tabular}{@{}ccccc@{}}
\toprule
\textbf{Prompt} & \textbf{FPGA (\iffour{})} & \textbf{CUDA (\bfteen{})} & \textbf{$\Delta$ steps} & \textbf{Agree} \\
\midrule
0 & \checkmark~94  & \checkmark~95  & $-1$  & \checkmark \\
1 & \checkmark~160 & \checkmark~126 & $+34$ & \checkmark \\
2 & \checkmark~123 & \checkmark~115 & $+8$  & \checkmark \\
3 & \checkmark~108 & \checkmark~109 & $-1$  & \checkmark \\
4 & $\times$~230 & $\times$~230 & --- & \checkmark \\
5 & $\times$~290 & $\times$~290 & --- & \checkmark \\
6 & \checkmark~121 & \checkmark~120 & $+1$  & \checkmark \\
7 & $\times$~230 & $\times$~230 & --- & \checkmark \\
8 & \checkmark~107 & \checkmark~106 & $+1$  & \checkmark \\
9 & $\times$~230 & $\times$~230 & --- & \checkmark \\
\midrule
\textbf{Rate} & \textbf{6/10} & \textbf{6/10} & & \textbf{10/10} \\
\bottomrule
\end{tabular}
\end{table}

\vspace{4pt}
\noindent\textbf{Takeaway.} \iffour{} is not just numerically close---it is
behaviorally identical. Across all 10 prompts the FPGA \iffour{} policy and the
\bfteen{} reference reach the \emph{same} success/failure outcome ($10/10$
agreement, $6/6$ shared successes), and on the successful episodes their
trajectory lengths differ by only a handful of steps ($\Delta$ within a few
percent, $+34$ on one longer reach). The $16.3$\,dB end-to-end SNR floor of
Sec.~\ref{sec:quant} sounds alarming in isolation, but the low-SNR components
are small-magnitude action dimensions that do not change what the robot
\emph{does}---the task-critical axes and gripper are preserved, and the
closed-loop behavior confirms it. Put as a benchmark number: quantization
\emph{did not regress task performance at all} --- same $6/10$ success rate,
same outcome on every prompt, successful trajectories within a few control
steps --- so the $\sim$8$\times$ checkpoint compression is free at the task
level on this suite.

%=============================================================================
\section{Throughput: FPGA vs.\ CUDA}
\label{sec:throughput}
%=============================================================================
We time each stage on both platforms and report it in the same form: per-section
wall-clock and the achieved rate (section FLOPs $\div$ time). FPGA times come
from the engine's own instrumentation. The
GPU baseline is openpi's JAX/XLA \texttt{Pi0} --- the deployment path, sustaining
$130.2$\,ms/inference ($7.68$\,Hz), batch~1 --- broken into the same three stages
by jitting each separately and forcing device completion
(\texttt{block\_until\_ready}) at the boundaries. Because these stages are
dispatched separately rather than fused into the deployed \texttt{lax.while\_loop},
their sum ($135.4$\,ms) slightly exceeds the fused end-to-end ($130.2$\,ms); the
per-stage split is used for the breakdown, the fused number for the headline.

\begin{table}[H]
\centering\small
\caption{RTX~5070 (\bfteen{}, openpi JAX/XLA) per-section timing and achieved
throughput. Batch~1, 10 denoise steps, sample~0. Per-stage times are
separately-jitted; End-to-end is the fused \texttt{policy.infer}.}
\label{tab:timing-gpu}
\begin{tabular}{@{}lcc@{}}
\toprule
\textbf{Model section} & \textbf{Time (ms)} & \textbf{GFLOP/s} \\
\midrule
Vision encoder          & 19.41  & 33\,793 \\
Prefix (Gemma-2B)       & 79.29  & 42\,877 \\
Action expert (denoise) & 36.70  & 2\,035  \\
\midrule
\textbf{End-to-end}     & 130.15 & 31\,736 \\
\textbf{Throughput (inf/s)} & \multicolumn{2}{c}{7.68} \\
\bottomrule
\end{tabular}
\end{table}

\begin{table}[H]
\centering\small
\caption{FPGA (\iffour{}) per-section timing and achieved throughput
(\ue{} v1 on Xilinx Kintex-7). Batch~1, 10 denoise steps, sample~0, 3-slot
vision, 832-token prefix. GFLOP/s and \%~peak are on the hardware-issued
(padded) FLOPs; theoretical peak is $26.6$\,GFLOP/s (one Apex Compute
Transformer Engine at $208$\,MHz).}
\label{tab:timing-fpga}
\setlength{\tabcolsep}{4pt}
\begin{tabular}{@{}lcccccc@{}}
\toprule
\textbf{Model section} & \textbf{Time (s)} & \textbf{Share} & \textbf{Eff.\ GFLOP} & \textbf{HW GFLOP} & \textbf{GFLOP/s (hw)} & \textbf{\% peak} \\
\midrule
Vision encoder          & 35.4  & 18.3\% & 655.9  & 850.7  & 24.0 & 90\% \\
Prefix (Gemma-2B)       & 137.6 & 71.1\% & 3399.8 & 3399.8 & 24.7 & 93\% \\
Action expert (denoise) & 20.6  & 10.6\% & 74.7   & 483.5  & 23.5 & 88\% \\
\midrule
\textbf{End-to-end}     & \textbf{193.6} & 100\% & \textbf{4130.4} & \textbf{4734.0} & \textbf{24.5} & \textbf{92\%} \\
\textbf{Throughput (inf/s)} & \multicolumn{6}{c}{0.0052} \\
\bottomrule
\end{tabular}
\end{table}

\begin{table}[H]
\centering\small
\caption{Silicon (\iffour{}, \emph{estimated}) per-section timing and achieved
throughput --- projected ASIC implementation of the \ue{} v1 IP at
$1.6$\,GHz with $100$\,GB/s DRAM, scaling the Apex Compute Transformer
Engine count in two chip variations ---
\textbf{Chip Variation~1} (8~engines, $\sim$5\,W) and \textbf{Chip
Variation~2} (16~engines, wattage TBD) --- $1638.4$ /
$3276.8$\,GFLOP/s aggregate peak.
Batch~1, 10 denoise steps. Linear $61.5\times$ / $123\times$ compute scaling
of the measured FPGA run at its achieved \%-of-peak; neither variation is
bandwidth-bound (see below).}
\label{tab:timing-silicon}
\begin{tabular}{@{}lcccc@{}}
\toprule
& \multicolumn{2}{c}{\textbf{Chip Variation 1}} & \multicolumn{2}{c}{\textbf{Chip Variation 2}} \\
& \multicolumn{2}{c}{\footnotesize (8 engines, $\sim$5\,W)} & \multicolumn{2}{c}{\footnotesize (16 engines, wattage TBD)} \\
\cmidrule(lr){2-3}\cmidrule(l){4-5}
\textbf{Model section} & \textbf{Time (ms)} & \textbf{GFLOP/s} & \textbf{Time (ms)} & \textbf{GFLOP/s} \\
\midrule
Vision encoder          & 575 & 1\,479 & 288 & 2\,958 \\
Prefix (Gemma-2B)       & 2\,236 & 1\,520 & 1\,118 & 3\,041 \\
Action expert (denoise) & 335 & 1\,444 & 167 & 2\,889 \\
\midrule
\textbf{End-to-end}     & 3\,146 & 1\,505 & 1\,573 & 3\,010 \\
\textbf{Throughput (inf/s)} & \multicolumn{2}{c}{0.32} & \multicolumn{2}{c}{0.64} \\
\bottomrule
\end{tabular}
\end{table}

\vspace{2pt}
\noindent\textbf{Why the projection is linear.} Going from one Apex Compute
Transformer Engine at
$208$\,MHz to 8 engines at $1.6$\,GHz multiplies compute by $61.5\times$
($26.6\!\to\!1638.4$\,GFLOP/s peak) but DRAM bandwidth by only $16.7\times$
($\sim$6\,$\to$\,$100$\,GB/s), so each section survives as compute-bound only
if its arithmetic intensity clears the ASIC ridge point of
$1638.4/100\approx16$\,FLOP/byte. Traffic here is dominated by streaming the
\iffour{} weights ($4.25$\,bits/weight, Sec.~\ref{sec:port}): the prefix
reads $\sim$1.1\,GB once
for its 3399.8\,GFLOP ($\sim$3200\,FLOP/byte), the vision tower $\sim$0.65\,GB
over three slots ($\sim$1300\,FLOP/byte), and the worst case --- the action
expert, whose $\sim$0.16\,GB are re-streamed on \emph{every} denoise step
against just $\sim$48\,GFLOP of padded work per step --- still lands at
$\gtrsim$200\,FLOP/byte with activation and KV-cache traffic included. All
three sections sit an order of magnitude above the ridge (worst-section
demand $<$10\,GB/s of the $100$ available), and doubling to 16 engines
(Chip Variation~2) only
lifts the ridge to $\approx$33\,FLOP/byte --- still $\geq$6$\times$ clear. So
neither variation is bandwidth-bound, and times extrapolate linearly at the
FPGA's measured $88$--$93\%$ of peak, assuming ideal engine scaling:
$3.15$\,s per inference ($0.32$\,inf/s) for Chip Variation~1 and $1.57$\,s
($0.64$\,inf/s) for Chip Variation~2 --- $61.5\times$ / $123\times$ the FPGA, within
$24\times$ / $12\times$ of the RTX~5070's $7.68$\,inf/s. Either way the
engine executes the whole vision--prefix--denoise pipeline autonomously:
host work per action chunk is an embedding gather on the way in and a
$10{\times}7$ readback on the way out, so the platform's CPU stays free for
the robot stack.

\vspace{2pt}
\noindent\textbf{Power.} The RTX~5070 is a $250$\,W-TGP board; Chip
Variation~1 is specified at $\sim$5\,W, and Chip Variation~2's wattage is
TBD. Table~\ref{tab:platforms} collects the platform comparison.

\begin{table}[H]
\centering\small
\caption{Platform comparison: measured RTX~5070 and \ue{}~v1 FPGA runs, and
the two projected chip variations. Achieved GFLOP/s counts effective FLOPs
on the GPU and hardware-issued FLOPs on the \ue{} platforms
(Tables~\ref{tab:timing-gpu}--\ref{tab:timing-silicon}); peaks are \bfteen{}
rates. The GPU peak is the whitepaper dense \bfteen{} tensor rate with FP32
accumulation; measured GPU rates above it reflect XLA matmuls dispatched with
reduced-precision accumulation, whose dense ceiling is $2\times$ higher.}
\label{tab:platforms}
\setlength{\tabcolsep}{4pt}
\begin{tabular}{@{}lcccc@{}}
\toprule
& \textbf{RTX 5070} & \textbf{FPGA (Kintex-7)} & \textbf{Chip Var.~1} & \textbf{Chip Var.~2} \\
\midrule
Compute            & 6144 CUDA cores & 1 engine & 8 engines & 16 engines \\
Clock              & $\sim$2.5\,GHz & 208\,MHz & 1.6\,GHz & 1.6\,GHz \\
Peak GFLOP/s (\bfteen{}) & 30\,870 & 26.6 & 1\,638 & 3\,277 \\
DRAM bandwidth     & 672\,GB/s & $\sim$6\,GB/s & 100\,GB/s & 100\,GB/s \\
Power              & 250\,W (TGP) & --- & $\sim$5\,W & TBD \\
Achieved GFLOP/s   & 31\,736 & 24.5 & 1\,505 & 3\,010 \\
Time / inference   & 0.130\,s & 193.6\,s & 3.15\,s & 1.57\,s \\
Throughput         & 7.68\,inf/s & 0.0052\,inf/s & 0.32\,inf/s & 0.64\,inf/s \\
\bottomrule
\end{tabular}
\end{table}

\vspace{2pt}
\noindent\textbf{Config.} FPGA: \ue{} v1 IP (one Apex Compute Transformer
Engine) on Xilinx Kintex-7 at $208$\,MHz
($26.6$\,GFLOP/s theoretical peak), \iffour{}
weights / \bfteen{} compute, $4$\,GB DRAM at $\sim$6\,GB/s. GPU: RTX~5070
(6144 CUDA cores, $250$\,W TGP), \bfteen{}, $12$\,GB GDDR7 at $672$\,GB/s,
openpi JAX/XLA. Silicon: Chip Variation~1 --- 8 engines at $1.6$\,GHz,
$\sim$5\,W; Chip Variation~2 --- 16 engines at $1.6$\,GHz, wattage TBD; both
$100$\,GB/s DRAM. All: batch~1, 10 denoise steps, sample~0.

\vspace{4pt}
The entire \pizero{} VLA --- vision, Gemma-2B prefix, and the 10-step
flow-matching expert --- runs on the \ue{}~v1. The recurring theme is the split
between where FLOPs concentrate (the compute-bound prefix) and where latency
does (the sequential, FLOP-light denoise loop), and how \iffour{} trades a
controlled, section-localized accuracy cost --- one with no measured
closed-loop task regression (Sec.~\ref{sec:libero}) --- for fitting a
$\sim$13\,GB model in a sub-4\,GB budget. Tables~\ref{tab:budget}--\ref{tab:platforms} complete
the characterization.

\end{document}
